\documentclass[10pt,a4paper]{amsart}
\usepackage[margin=19mm]{geometry}
\usepackage{amsmath,amssymb,mathtools}
\usepackage[scr=rsfs,cal=euler]{mathalfa}
\usepackage{xfrac}
\usepackage[unicode,hidelinks,bookmarks=false]{hyperref}
\newcommand{\ie}{i.e.\ }
\newcommand{\cf}{cf.\ }
\newcommand{\resp}{respectively\ }
\newcommand{\Subsection}{\subsection}
\providecommand{\nobreakdash}{\nobreak\hbox{-}}
\setlength{\emergencystretch}{2em}
\multlinegap0pt
\usepackage{amssymb}
\usepackage{mathtools}
\usepackage{xfrac}
\usepackage{tikz}
\usepackage[shortlabels]{enumitem}
\newtheorem{remark}{Remark}
\newtheorem{lemma}{Lemma}
\DeclareMathOperator{\Spec}{Spec}
\DeclareMathOperator{\TT}{\mathbb T}
\newcommand{\kk}{\Bbbk}
\title{Automorphism groups of semi-rigid $T$-varieties of complexity one}
\author{Kirill Rassolov}
\date{Source: arXiv:2609.09362v1 (CC BY 4.0)}
\begin{document}
\maketitle

\noindent\emph{Source and licence.} Automorphism groups of semi-rigid $T$-varieties of complexity one. Version arXiv:2609.09362v1 (CC BY 4.0), distributed by arXiv under the Creative Commons Attribution 4.0 International licence, http://creativecommons.org/licenses/by/4.0/, as recorded in the arXiv licence metadata for this exact version and shown on the abstract page https://arxiv.org/abs/2609.09362v1. Full licence notice: \texttt{licenses/arxiv-2609.09362-CC-BY-4.0.txt}.\par\smallskip

\noindent\textit{Editorial scope.} Counted unit: the article's lemma Trinomial-N-Lemma (source0.tex lines 269--275). The article's complete original proof of it is delivered with it (source0.tex lines 276--286, one \texttt{proof} environment of 11 lines). No mathematics is written, repaired or re-derived: the statement and the proof are the article's own lines, in the article's own notation, and no retained line is substituted or re-flowed.  Retained uncounted as the source-local context the proof needs: lines 252--254; lines 260--262.  Nothing was omitted from any retained span, so the package carries no omission record.  The retained lines cite 4 works, whose entries are transcribed verbatim from the article's own bibliography source in the e-print and delivered with short editorial labels recorded in the item's reference mapping.  No retained line contains a figure environment, an image-inclusion command or drawing code, so no image is delivered and the package declares no asset dependency.  Editorial formatting only, in addition: an AMS \texttt{amsart} preamble that restates the article's own declarations and macros used by the retained lines, namely the environments \texttt{remark}, \texttt{lemma} on separate counters, together with the standard packages those lines need.  The retained environments print in this document as Remark 1, Lemma 1 in order of appearance, whereas the article prints the same items with the numbering of its own full document.  The complete native source of the article as distributed in the arXiv e-print for this exact version is bundled unchanged as \texttt{sources/209808/source0.tex}.
\begin{remark} (see \cite{HW})\label{RemarkPoint}
    If $X(A)$ is of Type 2, then the action of the canonical torus $\TT$ is pointed, meaning that the rational polyhedral cone denerated by the degrees occurring in $R(A)$ contains no line.
\end{remark}

\begin{lemma}\label{LemmaHumphreys}
    Let $G$ be an ind-group containing a torus $T$ as a closed subgroup. Then $N_G(T)^\circ = C_G(T)^\circ$.
\end{lemma}

\begin{lemma}\label{Trinomial-N-Lemma}
    Let $X$ be a trinomial variety, $\TT$ the canonical torus of $X$.
    \begin{enumerate}
        \item[\emph{(i)}] In case of Type 1, $N(\TT)$ is a locally algebraic group having a torus as identity component.
        \item[\emph{(ii)}] In case of Type 2, $N(\TT)$ is a linear algebraic group.
    \end{enumerate} 
\end{lemma}

\begin{proof}
    Let $X$ be of Type $1$. %, and let $n_i$ stand for the number of the variables appearing in the monomial $T_i^{l_i}$. 
    If $n_i>1$ for some $i$, the assersion follows from~\cite[Lemma~6.2 and Corollary~3.8]{BorG}. %, , then %the variables $T_{ij}$, $j=1,\ldots n_i$, are isolated semi-invariants. Suppose there exists such $i$; then 
    %$\kk[X]$ is integral over the subalgebra generated by the isolated semi-invariants. It follows from
    %~\cite[Corollary~3.8]{BorG} that $N(\TT)^\circ$ is a torus.
    Thus it remains to treat the case $n_i=1$ for all $i$, i.e., each $T_{i}^{l_i}$ is a power of a variable. Then $X$ is either rigid or isomorphic to the affine space with coordinates $T_{11}, S_1,\ldots, S_m$ (see the rigidity criterion in~\cite{EGS}). In the first case, the assertion follows from~\cite[Proposition~6.3]{BorG}. 
    
    Now let $X\simeq \Spec \kk[T_{11}, S_1,\ldots, S_m]$ with $\TT\simeq (\kk^\times)^m$ acting on $S_1,\ldots, S_m$.  For any $\varphi\in N(\TT)$, the comorphism $\varphi^*$ leaves $\kk[T_{11}]$ invariant as functions of $\TT$-weight zero. Hence $\varphi^*T_{11}=aT_{11},$ $a\in\kk^\times$. By Lemma~\ref{LemmaHumphreys} we may further assume that $\varphi\in C(\TT)$. Thus $\varphi^*$ must preserve the weights of regular functions with respect to $\TT$, so $\varphi^*S_k=S_kf_k(T_{11})$ for all $k$, $f_k$ being polynomials in one variable. Since $\varphi^*$ is invertible, all the $f_k$ are constant. Therefore $N(\TT)^\circ= \kk^\times\times\TT$. 

    Finally, suppose that $X$ is of Type 2. Then the $\TT$-action is pointed (Remark~\ref{RemarkPoint}), hence $N(\TT)$ is a linear algebraic group by~\cite[Proposition~2.2]{AHHL}. 
\end{proof}

\begin{thebibliography}{99}
\bibitem[AHHL]{AHHL} Ivan Arzhantsev, J\"urgen Hausen, Elaine Herppich, and Alvaro Liendo. \emph{The automorphism group of a variety with torus action of complexity one.} Mosc. Math.~J.~14 (2014), no.~3, 429-471
\bibitem[BorG]{BorG} V. Borovik, S. Gaifullin. \emph{Isolated torus invariants and automorphism groups of rigid varieties}, Journal of Algebra 666, 821--839 (2025)
\bibitem[EGS]{EGS} Polina Evdokimova, Sergey Gaifullin, and Anton Shafarevich. {\it Rigid Trinomial Varieties.} arXiv:2307.06672v1 (2023).
\bibitem[HW]{HW} Hausen~J., Wrobel~M., \emph{Non-complete rational T-varieties of complexity one.} Math. Nachr.~{\bf 290} (2017), no.~5-6, 815--826.
\end{thebibliography}
\end{document}
